% RESEARCH_DRAFT_SHA256: aa0f37e31cca24a51c48b9cfe2ac92e5f1f0bb544aeba94e6ea57e0e843eaf87
% Layout and source-reference wrappers expanded for the web renderer.
\documentclass[11pt]{article}
\usepackage[margin=25mm]{geometry}
\usepackage{fontspec,unicode-math}
\setmainfont{Latin Modern Roman}
\setmathfont{Latin Modern Math}
\usepackage{amsmath,amssymb,mathtools,xurl,hyperref}
\hypersetup{colorlinks=true,urlcolor=blue,linkcolor=blue}
\setlength{\parindent}{0pt}
\setlength{\parskip}{5pt}
\setlength{\emergencystretch}{3em}
\newcommand{\sref}[2]{\hyperlink{source:#1:#2}{[S#2]}}
\newcommand{\cataloguescope}[1]{\paragraph{Source question.}#1}
\newcommand{\floor}[1]{\left\lfloor#1\right\rfloor}
\begin{document}

\section*{Stabilizer extent: a measurement bound and the failed amplification inference}
\noindent UnsolvedMath ID 30004861; OWR-8415352-003. Research attempt dated 11 October 2026.
\paragraph{Result in scope.}This manuscript contains a source-specific attempt and complete proofs of the stated auxiliary calculations. It makes no new solution or novelty claim for the entire target.
\section{Original problem and preserved definition}
The following definition and source list are reproduced from the repository. Any earlier inspection dates and reported literature claims in this preserved text describe that source record, not additional verification by this attempt.
% BEGIN PRESERVED DEFINITION
\subsection{Definitions and mathematical statement}

\paragraph{Shared notation.}Standard number systems are denoted by $\mathbb N,\mathbb Z,\mathbb Q,\mathbb R,\mathbb C$; index ranges, computational models and specialized notation follow the cited problem source and the definitions below.


For a normalized pure $n$-qubit state $\psi$, its stabilizer extent is $\xi(\psi)=\min\{(\sum_j|c_j|)^2:\psi=\sum_jc_j\phi_j,\ \phi_j\text{ normalized stabilizer states}\}$. A stabilizer state is a common eigenstate of a maximal commuting subgroup of the qubit Pauli group. Product decompositions give $\xi(\phi\otimes\psi)\leq\xi(\phi)\xi(\psi)$. The gap may be described additively or by the ratio $\xi(\phi)\xi(\psi)/\xi(\phi\otimes\psi)$.
\paragraph{Statement recorded in the supplied source.}
How large can the submultiplicativity gap between $\xi(\phi\otimes\psi)$ and $\xi(\phi)\xi(\psi)$ be for quantum states $\phi$ and $\psi$? \href{https://doi.org/10.4171/owr/2021/49}{[S2]}
\paragraph{Scope qualification.}
Extent is defined here for pure states, as in the source passage. The supplied question does not select an additive or multiplicative gap or fix the number of qubits; these must be stated in any quantitative bound.
\paragraph{Source status.}
The supplied catalogue labels this record partially solved. Its source-status metadata was checked against the supplied snapshot on 2026-10-11. This import does not independently establish current open status or mathematical correctness. The catalogue literature assessment is dated 2026-08-22.
\subsection{Short English statement}
Quantify how much entangled stabilizer decompositions can improve on taking the product of optimal decompositions.
\subsection{Sources}
\begin{itemize}
\item[S1] Ulam AI and the UnsolvedMath Contributors, UnsolvedMath: A Curated Collection of Open Mathematics Problems, record 30004861 (OWR-8415352-003). Catalogue curation and metadata: CC BY 4.0. The supplied source question is retained; definitions and the English summary are editorial additions.
\url{https://huggingface.co/datasets/ulamai/UnsolvedMath}.
\item[S2] Geometry and Optimization in Quantum Information (2022), Oberwolfach Reports, DOI 10.4171/owr/2021/49, pp. 2665--2720
\url{https://doi.org/10.4171/owr/2021/49}.
\textit{Underlying Oberwolfach Reports material retains its source-specific licence, including CC BY-SA 4.0 where applicable; see THIRD\_PARTY\_NOTICES.md.}
\end{itemize}
% END PRESERVED DEFINITION
\subsection{Verbatim supplied catalogue question}
How large can the submultiplicativity gap between $\xi(\phi\otimes\psi)$ and $\xi(\phi)\xi(\psi)$ be for quantum states $\phi$ and $\psi$?

\section{Source target and provenance of this attempt}
The primary report, printed pages 2680--2682, already establishes
strict submultiplicativity through the cited published paper and
then asks how large the gap can be. Its duplicated factor in one
printed line is interpreted as the product of the two extents,
consistent with the preceding definitions and the preserved
catalogue question. We treat normalized pure states and state
both multiplicative and additive upper bounds as functions of
the two qubit counts. These are auxiliary bounds with no
sharpness or novelty claim.
The root Codex GPT-6 model (exact serving revision not exposed)
suggested the measurement-branch proof in its local auxiliary
note. The drafting model audited that argument and supplies the
complete stabilizer-closure proof below. This is additional
model mathematical contribution, not human authorship or
independent human verification.

\section{Atomic norm and the measurement closure lemma}
For an arbitrary vector $v$ on $n$ qubits define
\[
 a_n(v)=\inf\left\{\sum_j|c_j|:v=\sum_j c_jS_j,
 \ S_j\text{ normalized stabilizer states}\right\}.
\]
The computational basis makes the feasible set nonempty.
The triangle inequality for the Euclidean norm and then
Cauchy--Schwarz in a computational-basis expansion give
\[
 \|v\|_2\le a_n(v)\le2^{n/2}\|v\|_2.
\]
Combining expansions proves the triangle inequality; scaling
coefficients and then scaling back proves
$a_n(zv)=|z|a_n(v)$. For normalized $v$, $a_n(v)^2=\xi(v)$.
Taking infima suffices throughout, so attainment is not needed.

We will use the following closure fact: measuring any subset of
qubits of a normalized stabilizer state in the computational
basis leaves, on every nonzero branch, a scalar times a normalized
stabilizer state on the other qubits. Here is a full algebraic
proof. An $N$-qubit stabilizer state has a group $G$ of $2^N$
commuting Hermitian Pauli operators fixing it and not containing
$-I$. The average $2^{-N}\sum_{g\in G}g$ is a projector: squaring
it counts each group element $2^N$ times, and it is Hermitian.
Its trace is one because only $I$ has nonzero trace. Thus its
range is the one-dimensional stabilized space. The same argument
shows that a commuting Pauli group of size $2^{N+1}$ with no
$-I$ on $N$ qubits is impossible, because its average would be
a projector of trace $1/2$.

Measure the Pauli $Z$ of one qubit. If $Z$ or $-Z$ belongs to
$G$, the outcome is deterministic. Otherwise some generator
anticommutes with $Z$: if all commuted, adjoining $Z$ would
give the impossible larger group just described. Choose such a
generator $g_1$. By multiplying any other anticommuting generator
by $g_1$, take the remaining $N-1$ independent generators to
commute with $Z$. For either outcome sign $z\in\{-1,1\}$,
replace $g_1$ by $zZ$. These $N$ generators are independent:
a relation involving $Z$ would place $\pm Z$ in $G$, and a
relation without it contradicts the original independence.
They commute. The projected vector is nonzero because
$\langle Z\rangle=\langle g_1Zg_1\rangle=-\langle Z\rangle=0$,
so each outcome has probability $1/2$, and the new generators
fix that vector. Thus they cannot generate $-I$.

In either case choose $zZ$ as one generator. Every other generator
commutes with $Z$ and hence has only $I$ or $Z$ on this qubit;
multiplying by $zZ$ removes any $Z$ there. Restricting these
$N-1$ independent generators to the remaining qubits preserves
independence and the absence of $-I$. Their unique stabilized
vector is the residual state. The measured qubit is fixed to its
computational basis outcome. Iterating proves the stated closure
lemma. For a normalized joint stabilizer $S_j$, write its branch
vectors as $\alpha_{jx}s_{jx}$; orthogonality of measurement
outcomes gives $\sum_x|\alpha_{jx}|^2=1$.

\section{A complete non-sharp bound on both gap conventions}
Let $\phi,\psi$ be normalized $n$- and $m$-qubit states. Product
expansions into stabilizers are stabilizer expansions, since the
tensor products of the two independent generator groups have
the required number of commuting independent generators.
Taking infima proves the usual upper bound. For the lower bound,
start with any joint expansion
\[
 \phi\otimes\psi=\sum_j c_jS_j.
\]
Measure the last $m$ qubits. The outcome $x$ gives
$\psi_x\phi=\sum_jc_j\alpha_{jx}s_{jx}$, so
\[
 |\psi_x|a_n(\phi)\le\sum_j|c_j||\alpha_{jx}|.
\]
Square and sum over outcomes. Normalization of $\psi$ and
Cauchy--Schwarz for each pair of branch-amplitude vectors imply
\[
 \xi(\phi)\le
 \sum_{i,j}|c_ic_j|\sum_x|\alpha_{ix}||\alpha_{jx}|
 \le\left(\sum_j|c_j|\right)^2.
\]
Infimize and exchange the factors to obtain
\[
 \max\{\xi(\phi),\xi(\psi)\}
 \le\xi(\phi\otimes\psi)\le\xi(\phi)\xi(\psi).
\]
The norm bounds above give $1\le\xi(\phi)\le2^n$ and
$1\le\xi(\psi)\le2^m$. Consequently the multiplicative ratio
and additive difference satisfy
\[
 1\le\frac{\xi(\phi)\xi(\psi)}{\xi(\phi\otimes\psi)}
 \le\min\{\xi(\phi),\xi(\psi)\}\le2^{\min(n,m)},
\]
\[
 0\le\xi(\phi)\xi(\psi)-\xi(\phi\otimes\psi)
 \le2^{n+m}-2^{\max(n,m)}.
\]
For the last inequality, $ab-\max(a,b)$ is nondecreasing in
each variable for $a,b\ge1$: on either side of $a=b$ its
slope in either variable is either the other variable or the other
variable minus one. Substitute $a\le2^n,b\le2^m$.
If one factor is a stabilizer, its extent is one by the one-term
expansion and the norm lower bound, so the two-sided inequality
proves exact equality of the product extent with the other extent.
The statements include a zero-qubit scalar factor, with extent one.

\section{Why a strict example does not amplify automatically}
A natural attempted route to a large gap takes tensor powers of
a known strictly submultiplicative pair. Write
\[
 A_t=\xi(\phi^{\otimes t}),\quad
 B_t=\xi(\psi^{\otimes t}),\quad
 C_t=\xi((\phi\otimes\psi)^{\otimes t}).
\]
Submultiplicativity gives $A_t\le A_1^t$, $B_t\le B_1^t$ and
$C_t\le C_1^t$. The numerator estimates are upper bounds, so
they do not imply $A_tB_t/C_t\ge(A_1B_1/C_1)^t$.
For an exact check of the logical failure, the abstract constant
sequences $A_t=B_t=2$, $C_t=3$ obey all these submultiplicative
inequalities, monotonicity in $t$, and
$\max(A_t,B_t)\le C_t\le A_tB_t$, but their ratio is always
$4/3$. These sequences are not claimed to be realized by states;
they show that the displayed inequalities alone do not amplify
a strict gap. Additional control of the factor tensor powers or
an explicit growing family is missing.

\section{Adversarial audit and remaining target}
The audit checked independence of replacement generators,
postselection normalization, the absolute values in the branch
Cauchy--Schwarz step, infima versus attained decompositions, and
the direction of all tensor-power inequalities. Arbitrary partial
inner products with magic states were not assumed to preserve
stabilizers. The known existence of strict inequality and the
source's low-qubit multiplicativity statements are context,
not re-proved or counted as new progress. No numerical experiment
or theorem prover was used. The sharp extremal or asymptotic
gap remains unresolved by this attempt.

\section{Process and contributions}
Mathematics and model-assisted adversarial self-audit: OpenAI gpt-6-astra, ultra reasoning, as dispatched for this research batch; exact serving revision not exposed. Codex (GPT-6; exact serving revision not exposed) materially assisted with orchestration and format review. No human mathematical authorship, mathematical edits, or independent human proof checking is asserted.
The coordinating Codex (GPT-6; exact serving revision not exposed) also provided a separate model-assisted written-proof review of this manuscript. This is not independent human verification.
Contribution prompt A was applied to the full definition and supplied source row, then prompt B was applied adversarially to the derived results. The item-specific audit above records the scope checks and rejected extrapolations. Local assembly and structural checks do not establish mathematical correctness. The source catalogue status is preserved in TOP\_PROBLEM and is separate from this attempt status.
\paragraph{Integration review.} Codex (GPT-6; exact serving revision not exposed) reviewed the written arguments and expanded layout and source-reference wrappers. Repository metadata and the production web renderer were checked using Python 3.13 and JavaScript with MathJax 3.2.2. This model-assisted review is not independent human proof verification; no PDF compilation or formal proof checking is asserted.

\section{Completion Estimate}
\noindent\textbf{COMPLETION: 0\%}
The proved material is a rederivation of known or auxiliary facts rather than a substantiated new advance on the original target. No sharp upper bound, matching family, or asymptotic extremal growth is obtained. The measurement bound is a known-style auxiliary estimate; tensor-power amplification lacks lower bounds for factor extents.
\section{References}
The complete original Sources subsection above is retained, including catalogue attribution and licence notices. The additional list records primary sources actually consulted for this attempt on 11 October 2026; uncited source links reproduced above were not all rechecked.
\begin{itemize}
\item OWR49/2021, Geometry and Optimization in Quantum Information; primary source passages read from cached report, with precise pages identified in each manuscript. \url{https://doi.org/10.4171/owr/2021/49}.
\item Heimendahl, Montealegre-Mora, Vallentin and Gross, Stabilizer extent is not multiplicative, Quantum5,400(2021); primary publication abstract inspected for established context, not used as an external proof dependency. \url{https://doi.org/10.22331/q-2021-02-24-400}.
\end{itemize}
\section{Final status}
\textbf{Unresolved in this attempt.} No sharp upper bound, matching family, or asymptotic extremal growth is obtained. The measurement bound is a known-style auxiliary estimate; tensor-power amplification lacks lower bounds for factor extents.
\end{document}
